\documentclass[10pt,a4paper]{amsart}
\usepackage[margin=19mm]{geometry}
\usepackage{amsmath,amssymb,mathtools}
\usepackage[scr=rsfs,cal=euler]{mathalfa}
\usepackage{xfrac}
\usepackage[unicode,hidelinks,bookmarks=false]{hyperref}
\newcommand{\ie}{i.e.\ }
\newcommand{\cf}{cf.\ }
\newcommand{\resp}{respectively\ }
\newcommand{\Subsection}{\subsection}
\providecommand{\nobreakdash}{\nobreak\hbox{-}}
\setlength{\emergencystretch}{2em}
\multlinegap0pt
\usepackage{mathtools,mathrsfs}
\usepackage{booktabs,array}
\numberwithin{equation}{section}
\newtheorem{theorem}{Theorem}[section]
\newtheorem{proposition}[theorem]{Proposition}
\newtheorem{lemma}[theorem]{Lemma}
\newtheorem{corollary}[theorem]{Corollary}
\newtheorem*{maintheorem}{Main theorem}
\newtheorem{definition}[theorem]{Definition}
\newtheorem{example}[theorem]{Example}
\newtheorem{assumption}[theorem]{Assumption}
\newtheorem{remark}[theorem]{Remark}
\newcommand{\C}{\mathbb C}
\newcommand{\R}{\mathbb R}
\newcommand{\Z}{\mathbb Z}
\newcommand{\T}{\mathcal T}
\newcommand{\Tgr}{\mathcal T_{g,r}}
\newcommand{\MB}{\mathcal M_B}
\newcommand{\MDol}{\mathcal M_{\mathrm{Dol}}}
\newcommand{\Mpar}{\mathcal M^{\mathrm{par}}_{\mathrm{Dol}}}
\newcommand{\ad}{\operatorname{ad}}
\newcommand{\End}{\operatorname{End}}
\newcommand{\Hom}{\operatorname{Hom}}
\newcommand{\Sym}{\operatorname{Sym}}
\newcommand{\tr}{\operatorname{tr}}
\newcommand{\Pf}{\operatorname{Pf}}
\newcommand{\Res}{\operatorname{Res}}
\newcommand{\rk}{\operatorname{rk}}
\newcommand{\SL}{\operatorname{SL}}
\newcommand{\SO}{\operatorname{SO}}
\newcommand{\Sp}{\operatorname{Sp}}
\newcommand{\SU}{\operatorname{SU}}
\newcommand{\pardeg}{\operatorname{pardeg}}
\newcommand{\Ad}{\operatorname{Ad}}
\begin{document}
\title[Quadratic one-forms on logarithmic Higgs moduli]{Quadratic one-forms on logarithmic Higgs moduli}
\author{Sumit Roy}
\date{}
\maketitle
\noindent\emph{Source and licence.} Quadratic one-forms on logarithmic Higgs moduli. Version arXiv:2606.15560v2 (CC BY 4.0). This material is distributed under the Creative Commons Attribution 4.0 International licence, http://creativecommons.org/licenses/by/4.0/, as recorded in the arXiv OAI licence metadata for this exact version and shown on the abstract page https://arxiv.org/abs/2606.15560v2. Full rights notice: licenses/arxiv-2606.15560-CC-BY-4.0.txt.\par\smallskip
\noindent\textit{Editorial scope.} Counted unit: the original \texttt{theorem} environment \texttt{thm:log-oneform} at source lines 2019--2043 together with its complete proof at lines 2045--2065; the delivered document keeps source lines 1200--2066 so that every referenced label and every citation resolves inside it. Native editable source is the paper's own concatenated TeX; the AMS wrapper is editorial formatting only. No publisher class, upstream script or unrelated artwork is redistributed.
\label{sec:prelim}

Let
\(
        D=p_1+\cdots+p_r
\)
be a reduced divisor on \(C\), and put
\(
        C^\circ=C\setminus D.
\)
A logarithmic Higgs field on \((C,D)\) is allowed simple poles along \(D\). In a
coordinate \(z\) centred at \(p_i\), it has the form
\[
        \Phi=
        \left(
        \frac{N_i}{z}+A_i(z)
        \right)dz,
        \qquad
        N_i=\operatorname{Res}_{p_i}(\Phi).
\]
The nilpotency of \(N_i\) controls the leading pole terms of invariant
polynomials in \(\Phi\).


\subsection{Parabolic Higgs bundles}

\begin{definition}
A \textit{parabolic bundle} over \((C,D)\) is a holomorphic vector bundle \(E\)
on \(C\), together with, for every \(p_i\in D\), a decreasing filtration
\[
        E_{p_i}=F_{i,1}\supset F_{i,2}\supset\cdots
        \supset F_{i,\ell_i}\supset F_{i,\ell_i+1}=0
\]
and real numbers
\[
        0\leq \alpha_{i,1}<\alpha_{i,2}<\cdots<\alpha_{i,\ell_i}<1.
\]
The numbers \(\alpha_{i,j}\) are called the \textit{parabolic weights}.
\end{definition}

The multiplicity of the weight \(\alpha_{i,j}\) is
\[
        m_{i,j}
        =
        \dim(F_{i,j}/F_{i,j+1}).
\]
The parabolic degree of \(E\) is
\[
        \operatorname{pardeg}(E)
        =
        \deg(E)
        +
        \sum_{i=1}^{r}\sum_{j=1}^{\ell_i}
        \alpha_{i,j}m_{i,j},
\]
and the parabolic slope is
\[
        \operatorname{par}\mu(E)
        =
        \frac{\operatorname{pardeg}(E)}{\operatorname{rk}(E)}.
\]

If \(F\subset E\) is a holomorphic subbundle, then \(F\) has an induced
parabolic structure.  At \(p_i\), the filtration is obtained by intersecting
with the filtration of \(E_{p_i}\)
\[
        F_{p_i}\cap F_{i,1}\supset
        F_{p_i}\cap F_{i,2}\supset\cdots\supset
        F_{p_i}\cap F_{i,\ell_i}\supset0,
\]
after deleting repetitions. The weights are the corresponding weights of
\(E\). The parabolic degree and parabolic slope of \(F\) are computed with
this induced parabolic structure.

\begin{definition}
A logarithmic Higgs field on \(E\) is a section
\(
        \Phi\in H^0(C,\operatorname{End}(E)\otimes K_C(D)).
\)
A \textit{parabolic Higgs bundle} is a parabolic vector bundle \(E\) together with a
logarithmic Higgs field \(\Phi\) such that, for every \(p_i\in D\),
\[
        \operatorname{Res}_{p_i}(\Phi)(F_{i,j})
        \subset F_{i,j}
        \qquad
        \text{for all }j.
\]
It is called \textit{strongly parabolic} if
\[
        \operatorname{Res}_{p_i}(\Phi)(F_{i,j})
        \subset F_{i,j+1}
        \qquad
        \text{for all }i,j.
\]
\end{definition}

Thus, for a strongly parabolic Higgs field, the residue at \(p_i\) is nilpotent
with respect to the flag at \(p_i\). In particular, if the flag is full, then
the residue is represented by a strictly upper triangular matrix in a basis
adapted to the flag.

\begin{definition}
A parabolic subbundle \(F\subset E\) is called \(\Phi\)-invariant if
\(
        \Phi(F)\subset F\otimes K_C(D).
\)
A parabolic Higgs bundle \((E,\Phi)\) is called \textit{stable} (resp. \textit{semistable}), if for
every proper nonzero \(\Phi\)-invariant parabolic subbundle \(F\subset E\), we have
\[
        \operatorname{par}\mu(F)<\operatorname{par}\mu(E),
        \quad (\text{resp.} \leq)        
\]
\end{definition}

For fixed rank, degree, parabolic weights and multiplicities, the moduli space
of semistable parabolic Higgs bundles is a quasi-projective variety.  We shall
only use its smooth stable locus (see \cite{MehtaSeshadri, Yokogawa, BodenYokogawa, BiswasParabolic} for more details). For the relation between parabolic Higgs bundles and Teichmuller
spaces of punctured surfaces, see \cite{BiswasGastesiGovindarajan}.

\begin{remark}
For \(SL(2,\mathbb C)\), we have
\(
        \det E\simeq\mathcal O_C,
        \ \text{and} \
        \Phi\in H^0(C,\operatorname{End}_0(E)\otimes K_C(D)).
\)
If \((E,\Phi)\) is strongly parabolic and
\(
        N_i=\operatorname{Res}_{p_i}(\Phi),
\)
then
\(
        N_i\in\mathfrak{sl}_2(\mathbb C).
\)
Since \(N_i\) strictly lowers the parabolic flag, it is nilpotent, in rank
two this gives
\(
        N_i^2=0,
\)
or equivalently
\(
        \operatorname{tr}(N_i^2)=0.
\)
This observation will be used later for
\(\operatorname{tr}(\Phi^2)\).
\end{remark}


\subsection{Logarithmic principal Higgs bundles}

\begin{definition}
A \textit{logarithmic \(G\)-Higgs bundle} on \((C,D)\) is a pair
\(
        (P,\Phi),
\)
where \(P\) is a holomorphic principal \(G\)-bundle on \(C\), and
\(
        \Phi\in H^0(C,\operatorname{ad}(P)\otimes K_C(D)).
\)
\end{definition}

Let \(p_i\in D\). Choose a coordinate \(z\) centred at \(p_i\), and choose a
local trivialization of \(\operatorname{ad}(P)\). Then the Higgs field has the
local form
\[
        \Phi=
        \left(
        \frac{N_i}{z}+A_i(z)
        \right)dz,
\]
where \(A_i(z)\) is holomorphic. The element
\(
        N_i\in\mathfrak g
\)
is the residue of \(\Phi\) at \(p_i\), and we write
\(
        N_i=\operatorname{Res}_{p_i}(\Phi).
\)

\begin{definition}
We say that the logarithmic Higgs field \(\Phi\) has nilpotent residues if
\(
        \operatorname{Res}_{p_i}(\Phi)
\)
is nilpotent for every \(p_i\in D\).
\end{definition}

For a reductive Lie algebra
\(
        \mathfrak g
        =
        \mathfrak z(\mathfrak g)\oplus[\mathfrak g,\mathfrak g],
\)
where \(\mathfrak z(\mathfrak g)\) is the centre and
\([\mathfrak g,\mathfrak g]\) is semisimple, an element
\(
        N\in\mathfrak g
\)
is called nilpotent if its component in \(\mathfrak z(\mathfrak g)\) is zero
and its component in \([\mathfrak g,\mathfrak g]\) is nilpotent.

We fix a smooth stable family
\[
        \mathcal M^{\mathrm{nil}}_{\mathrm{Dol}}(G)
        \longrightarrow
        \mathcal T_{g,r},
\]
whose fibre over \((C,D)\) is a smooth stable locus in the logarithmic Dolbeault
moduli space of \(G\)-Higgs bundles with fixed topological type and nilpotent
residues.

\subsection{Pointed Teichmuller space}

Let \(\mathcal T_{g,r}\) denote the Teichmuller space of pointed Riemann
surfaces of type \((g,r)\).  At the point represented by \((C,D)\), one has
\[
        T_{(C,D)}\mathcal T_{g,r}
        \cong
        H^1(C,T_C(-D)).
\]
Since
\[
        (T_C(-D))^\vee=K_C(D),
\]
Serre duality gives
\[
        T^*_{(C,D)}\mathcal T_{g,r}
        \cong
        H^0(C,K_C^2(D)).
\]
Thus the cotangent space is the space of meromorphic quadratic differentials
with at most simple poles at the marked points.
This is the standard cotangent-space description of pointed Teichmuller space;
see, for example, \cite{Gardiner,Hubbard}.

We shall use two representatives of a tangent vector
\(
        Y\in H^1(C,T_C(-D)).
\)
The first is a Dolbeault representative
\(
        \nu\in A^{0,1}(C,T_C(-D)),
\)
and the second is a \v{C}ech--Dolbeault representative
\(
        (\mu,\xi_1,\ldots,\xi_r),
\)
where \(\mu\) is a Beltrami differential on \(C^\circ\), and \(\xi_i\) is a
local holomorphic vector field near \(p_i\).  The vector
\(
        \xi_i(p_i)\in T_{p_i}C
\)
is the tangent direction of the marked point \(p_i\).

The following lemma will be used in the proof of the logarithmic energy-variation
formula.

\begin{lemma}
\label{lem:vanishing-representative}
Every class
\(
        Y\in H^1(C,T_C(-D))
\)
admits a Dolbeault representative
\(
        \nu\in A^{0,1}(C,T_C(-D))
\)
which is identically zero on a neighbourhood of \(D\).
\end{lemma}

\begin{proof}
Let
\(
        \nu_0\in A^{0,1}(C,T_C(-D))
\)
be a \(\bar\partial\)-closed representative of \(Y\). Choose pairwise disjoint
coordinate discs \(U_i\) around the points \(p_i\). Since \(U_i\) is
contractible, the local \(\bar\partial\)-Poincare lemma for the line bundle
\(T_C(-D)\) gives
\(
        \eta_i\in A^{0,0}(U_i,T_C(-D))
\)
such that
\(
        \bar\partial\eta_i=\nu_0|_{U_i}.
\)
Choose a smooth function \(\chi_i\) on \(C\), supported in \(U_i\), such that
\[
        \chi_i=1
        \quad\text{on a smaller disc }U_i'\subset U_i.
\] Define
\[
        \nu
        =
        \nu_0-\sum_{i=1}^{r}\bar\partial(\chi_i\eta_i).
\]
Then \(\nu\) differs from \(\nu_0\) by a \(\bar\partial\)-exact term, and hence
represents the same class in \(H^1(C,T_C(-D))\).  On \(U_i'\), we have
\(\chi_i=1\), and therefore
\[
        \nu
        =
        \nu_0-\bar\partial\eta_i
        =
        0.
\]
Thus \(\nu\) vanishes on a neighbourhood of \(D\).
\end{proof}

\subsection{Residue and the local Serre pairing}

For a meromorphic one-form \(\alpha\) near \(p\), we write
\(
        \operatorname{Res}_{p}(\alpha)
\)
for the coefficient of \(dz/z\) in a local coordinate \(z\) centred at \(p\). Equivalently,
\[
        \operatorname{Res}_{p}(\alpha)
        =
        \frac{1}{2\pi i}
        \int_{|z|=\varepsilon}\alpha .
\]
Let
\(
        q\in H^0(C,K_C^2(D)),
\)
and let
\(
        Y\in H^1(C,T_C(-D))
\)
be represented by
\(
        (\mu,\xi_1,\ldots,\xi_r),
\)
where \(\mu\) is a Beltrami differential on \(C^\circ\), and each \(\xi_i\) is
a local holomorphic vector field near \(p_i\). Then the Serre pairing is
\begin{equation}
        \langle q,Y\rangle
        =
        \int_{C^\circ}q\,\mu
        +
        \sum_{i=1}^{r}
        \operatorname{Res}_{p_i}(q\,\xi_i).
        \label{eq:serre-local}
\end{equation}
This is the \v{C}ech--Dolbeault form of the pairing
\[
        H^0(C,K_C^2(D))\times H^1(C,T_C(-D))
        \longrightarrow
        \mathbb C.
\]
This is the usual \v{C}ech--Dolbeault expression for the Serre pairing; see,
for example, \cite[Ch.~2]{GriffithsHarris}.
The residue term depends only on the tangent vector
\(
        \xi_i(p_i)\in T_{p_i}C.
\)
Indeed, if \(\eta_i\) is a local holomorphic vector field with
\(
        \eta_i(p_i)=0,
\)
then \(q\,\eta_i\) is holomorphic at \(p_i\), since \(q\) has at most a simple
pole. Therefore
\[
        \operatorname{Res}_{p_i}(q\,\eta_i)=0.
\]

\section{Nilpotent residues and invariant polynomials}
\label{sec:invpoly}

We use the notion of nilpotent residue fixed in Section~\ref{sec:prelim}. The
following lemma is about the quadratic invariant associated to an invariant
symmetric bilinear form.

We shall use the following consequence of the Jacobson-Morozov
theorem: if \(N\) is a nilpotent element of a complex reductive Lie algebra in
the sense fixed above, then there exists
\(
        H\in[\mathfrak g,\mathfrak g]
\)
such that
\(
        [H,N]=2N.
\)
(see \cite[Ch.~3]{CollingwoodMcGovern}).

\begin{proposition}
\label{prop:nilpotent-isotropic}
Let
\(
        B:\mathfrak g\otimes\mathfrak g\longrightarrow\mathbb C
\)
be an invariant symmetric bilinear form. If \(N\in\mathfrak g\) is nilpotent,
then
\(
        B(N,N)=0.
\)
\end{proposition}

\begin{proof}
By the definition fixed in Section~\ref{sec:prelim}, the central component of
\(N\) is zero. Hence
\(
        N\in[\mathfrak g,\mathfrak g].
\)
Choose
\(
        H\in[\mathfrak g,\mathfrak g]
\)
such that
\(
        [H,N]=2N.
\)
Using the invariance of \(B\), we get
\[
        2B(N,N)
        =
        B([H,N],N)
        =
        B(H,[N,N])
        =
        0.
\]
Therefore
\(
        B(N,N)=0. \qedhere
\)
\end{proof}

We also need the corresponding vanishing for homogeneous invariant polynomials.

\begin{theorem}
\label{thm:invpoly}
Let
\(
        p\in\mathbb C[\mathfrak g]^G
\)
be a homogeneous invariant polynomial of degree \(d\geq 1\). Let
\((P,\Phi)\) be a logarithmic \(G\)-Higgs bundle on \((C,D)\). Assume that
every residue
\(
        N_i=\operatorname{Res}_{p_i}(\Phi)
\)
is nilpotent. Then
\[
        p(\Phi)\in H^0\bigl(C,K_C^d((d-1)D)\bigr).
\]
\end{theorem}

\begin{proof}
Fix \(p_i\in D\).  Choose a coordinate \(z\) centred at \(p_i\), and choose a
local trivialization of \(\operatorname{ad}(P)\).  Then
\[
        \Phi=
        \left(
        \frac{N_i}{z}+A_i(z)
        \right)dz,
\]
where \(A_i(z)\) is holomorphic and
\(
        N_i=\operatorname{Res}_{p_i}(\Phi).
\)
Since \(p\) is homogeneous of degree \(d\),
\[
        p(\Phi)
        =
        z^{-d}p(N_i+zA_i(z))\,dz^d.
\]
We first show that
\(
        p(N_i)=0.
\)
By the nilpotency of \(N_i\) and the Jacobson-Morozov consequence, there exists
\(
        H_i\in[\mathfrak g,\mathfrak g]
\)
such that
\(
        [H_i,N_i]=2N_i.
\)
Since \(p\) is invariant under the adjoint action,
\[
        p(N_i)
        =
        p(e^{s\operatorname{ad}H_i}N_i)
        =
        p(e^{2s}N_i)
        =
        e^{2ds}p(N_i)
\]
for every \(s\in\mathbb C\).  Since \(d>0\), this forces
\(
        p(N_i)=0,
\)
which is the constant term of
\(
        p(N_i+zA_i(z)).
\)
Hence
\[
        p(N_i+zA_i(z))=O(z),
\]
and therefore
\[
        p(\Phi)=O(z^{-(d-1)})\,dz^d.
\]
Thus \(p(\Phi)\) has pole order at most \(d-1\) at \(p_i\).  Since this holds
for every \(p_i\in D\), we obtain
\[
        p(\Phi)\in H^0\bigl(C,K_C^d((d-1)D)\bigr).
\]
\end{proof}

We shall also use the following local refinement of the same argument.

\begin{lemma}
\label{lem:refined-pole}
Let \(p\in\mathbb C[\mathfrak g]^G\) be homogeneous of degree \(d\). Fix
\(p_i\in D\). Suppose that \(1\leq r_i\leq d\) and 
\[
        (d^j p)_{N_i}=0
        \qquad
        \text{for }0\leq j\leq r_i-1.
\]
Then \(p(\Phi)\) has pole order at most \(d-r_i\) at \(p_i\).
\end{lemma}

\begin{proof}
With the notation of the above proof,
\[
        p(\Phi)
        =
        z^{-d}p(N_i+zA_i(z))\,dz^d.
\]
Then the vanishing of the first \(r_i\) Taylor terms gives
\[
        p(N_i+zA_i(z))=O(z^{r_i}).
\]
Hence
\[
        p(\Phi)=O(z^{-(d-r_i)})\,dz^d.
\]
\end{proof}


\begin{corollary}
\label{cor:nilpotent-hitchin-map}
Let
\(
        p_1,\ldots,p_\ell
\)
be homogeneous generators of \(\mathbb C[\mathfrak g]^G\), with degrees
\(
        d_1,\ldots,d_\ell.
\)
On the nilpotent-residue locus, the logarithmic Hitchin map takes values in
\[
        \bigoplus_{j=1}^{\ell}
        H^0\bigl(C,K_C^{d_j}((d_j-1)D)\bigr).
\]
Equivalently, the usual logarithmic Hitchin map factors as
\[
        h_{G,\log}^{\mathrm{nil}}:
        \mathcal M_{\mathrm{Dol}}^{\mathrm{nil}}(G)
        \longrightarrow
        \bigoplus_{j=1}^{\ell}
        H^0\bigl(C,K_C^{d_j}((d_j-1)D)\bigr).
\]
\end{corollary}


\begin{proof}
This follows by applying Theorem~\ref{thm:invpoly} to each generator
\(p_j\).
\end{proof}




\begin{corollary}
\label{cor:quadratic}
Let
\(
        B:\mathfrak g\otimes\mathfrak g\longrightarrow\mathbb C
\)
be an invariant symmetric bilinear form.  Let \((P,\Phi)\) be a logarithmic
\(G\)-Higgs bundle on \((C,D)\) with nilpotent residues. Then
\[
        B(\Phi,\Phi)\in H^0(C,K_C^2(D)).
\]
\end{corollary}

\begin{proof}
The polynomial
\[
        p(X)=\frac12 B(X,X)
\]
is invariant and homogeneous of degree two.  The result follows from
Theorem~\ref{thm:invpoly}.
\end{proof}


\subsection{Sharpness for trace invariants}

For \(G=GL(n,\mathbb C)\), Theorem~\ref{thm:invpoly} applied to
\(
        p(X)=\operatorname{tr}(X^k)
\)
gives
\[
        \operatorname{tr}(\Phi^k)
        \in
        H^0\bigl(C,K_C^k((k-1)D)\bigr)
\]
whenever the residues of \(\Phi\) are nilpotent.  The following local
calculation shows that the pole order \(k-1\) cannot be improved in general.

\begin{proposition}
\label{prop:trace-sharp}
Let \(G=GL(n,\mathbb C)\), and let \(2\leq k\leq n\).  There are local logarithmic Higgs fields with nilpotent residue for which
\(
        \operatorname{tr}(\Phi^k)
\)
has a pole of order exactly \(k-1\).
\end{proposition}

\begin{proof}
It is enough to give a local logarithmic Higgs field with nilpotent residue for
which the coefficient of the pole of order \(k-1\) is nonzero.

Let
\[
        N=E_{12}+E_{23}+\cdots+E_{k-1,k}
        \in\mathfrak{gl}_k\subset\mathfrak{gl}_n
\]
be the nilpotent Jordan block of size \(k\). Choose a holomorphic matrix
\(A(z)\) with
\(
        A(0)=E_{k1}.
\)
Consider locally
\[
        \Phi=
        \left(
        \frac{N}{z}+A(z)
        \right)dz .
\]
The coefficient of \(z^{-(k-1)}dz^k\) in \(\operatorname{tr}(\Phi^k)\) is
obtained by taking exactly one factor \(A(0)\) and \(k-1\) factors \(N\). Hence
it is
\[
        \sum_{j=0}^{k-1}
        \operatorname{tr}
        \bigl(
        N^jA(0)N^{k-1-j}
        \bigr).
\]
By cyclicity of the trace, this equals
\[
        k\,\operatorname{tr}\bigl(N^{k-1}A(0)\bigr).
\]
Since
\[
        N^{k-1}=E_{1k},
        \qquad
        A(0)=E_{k1},
\]
we get
\[
        \operatorname{tr}\bigl(N^{k-1}A(0)\bigr)
        =
        \operatorname{tr}(E_{11})
        =
        1.
\]
Thus the coefficient is \(k\neq0\).  Therefore a pole of order \(k-1\) can occur.
\end{proof}

\begin{remark}
For \(k=2\), the nilpotent residue removes the double pole of
\(\operatorname{tr}(\Phi^2)\), but the simple pole may remain.  This is exactly
the local behaviour used later for strongly parabolic logarithmic
\(SL(2,\mathbb C)\)-Higgs fields.
\end{remark}

\subsection{The Pfaffian invariant}

Let
\(
        G=SO(2m,\mathbb C).
\)
The Pfaffian is the degree \(m\) invariant polynomial
\[
        \operatorname{Pf}:\mathfrak{so}(2m,\mathbb C)\longrightarrow\mathbb C.
\]
For a logarithmic \(SO(2m,\mathbb C)\)-Higgs field, the invariant
\(
        \operatorname{Pf}(\Phi)
\)
is therefore a meromorphic \(m\)-differential.  Theorem~\ref{thm:invpoly}
gives the pole estimate below. Under an additional rank condition on the
residues, the first Taylor term of the Pfaffian also vanishes, and the pole
order drops by one more.

\begin{corollary}
\label{cor:pfaffian-log}
Let \((P,\Phi)\) be a logarithmic \(SO(2m,\mathbb C)\)-Higgs bundle with
nilpotent residues. Then
\[
        \operatorname{Pf}(\Phi)
        \in
        H^0\bigl(C,K_C^m((m-1)D)\bigr).
\]
If, for every \(p_i\in D\),
\[
        \operatorname{rk}\operatorname{Res}_{p_i}(\Phi)\leq 2m-4,
\]
then
\[
        \operatorname{Pf}(\Phi)
        \in
        H^0\bigl(C,K_C^m((m-2)D)\bigr).
\]
\end{corollary}

\begin{proof}
The first statement follows from Theorem~\ref{thm:invpoly}, since
\(\operatorname{Pf}\) is homogeneous of degree \(m\). For the second assertion, put
\(
        N_i=\operatorname{Res}_{p_i}(\Phi).
\)
It is enough, by Lemma~\ref{lem:refined-pole}, to prove
\[
        d(\operatorname{Pf})_{N_i}=0.
\]
Choose a local orthogonal trivialization near \(p_i\), so that
\(N_i\) is represented by a skew-symmetric \(2m\times 2m\) matrix. For a skew-symmetric \(2m\times 2m\) matrix \(X\), the first partial
derivatives of \(\operatorname{Pf}(X)\) are, up to sign, Pfaffians of the
\((2m-2)\times(2m-2)\) skew-symmetric matrices obtained by deleting two
corresponding rows and columns. If
\[
        \operatorname{rk}(N_i)\leq 2m-4,
\]
then each such \((2m-2)\times(2m-2)\) matrix is singular. Its determinant is
zero, and hence its Pfaffian is zero. Therefore
\[
        d(\operatorname{Pf})_{N_i}=0.
\]
Then by Lemma~\ref{lem:refined-pole}, with \(p=\operatorname{Pf}\), \(d=m\), and
\(r_i=2\), we get the required bound.
\end{proof}

\begin{remark}
The quadratic invariant and the Pfaffian enter the logarithmic theory in
different degrees. The quadratic invariant gives
\[
        B(\Phi,\Phi)\in H^0(C,K_C^2(D)),
\]
which is the cotangent space to pointed Teichmuller space. The Pfaffian gives
\[
        \operatorname{Pf}(\Phi)\in H^0(C,K_C^m((m-1)D)),
\]
and, under the rank condition above,
\[
        \operatorname{Pf}(\Phi)\in H^0(C,K_C^m((m-2)D)).
\]
Thus, for \(m>2\), the Pfaffian belongs to the higher Hitchin Hamiltonians
rather than to the ordinary Teichmuller cotangent direction.  In the exceptional
case \(m=2\), the Pfaffian also has degree two.
\end{remark}

\section{The logarithmic quadratic one-form}
\label{sec:logoneform}

Let
\[
        \pi:\mathcal M^{\mathrm{nil}}_{\mathrm{Dol}}(G)
        \longrightarrow
        \mathcal T_{g,r}
\]
be the smooth stable family fixed in Section~\ref{sec:prelim}. Thus the fibre
over \((C,D)\) consists of logarithmic \(G\)-Higgs bundles
\(
        (P,\Phi),
        \ \text{where} \
        \Phi\in H^0(C,\operatorname{ad}(P)\otimes K_C(D)),
\)
with nilpotent residues. By Corollary~\ref{cor:quadratic},
\[
        B(\Phi,\Phi)\in H^0(C,K_C^2(D)).
\]
On the other hand,
\[
        T^*_{(C,D)}\mathcal T_{g,r}
        \cong
        H^0(C,K_C^2(D)).
\]
Hence the quadratic differential
\(
        -\frac12B(\Phi,\Phi)
\)
is a cotangent vector to pointed Teichmuller space at \((C,D)\).  We denote the
corresponding one-form on the family by
\(
        \phi_{G,\log}.
\)
Thus, for
\[
        Y\in T_{(C,D)}\mathcal T_{g,r}
        \cong H^1(C,T_C(-D)),
\]
we set
\[
        \phi_{G,\log}(Y)
        =
        -\frac12
        \left\langle B(\Phi,\Phi),Y\right\rangle .
\]


\begin{theorem}
\label{thm:log-oneform}
The one-form
\(
        \phi_{G,\log}
\)
is a holomorphic section of
\(
        \pi^*\Omega^1_{\mathcal T_{g,r}}
\)
over
\(
        \mathcal M^{\mathrm{nil}}_{\mathrm{Dol}}(G).
\)
If \(Y\) is represented by a Dolbeault form
\[
        \nu\in A^{0,1}(C,T_C(-D))
\]
which vanishes near \(D\), then
\[
        \phi_{G,\log}(Y)
        =
        -\frac12\int_C B(\Phi,\Phi)\nu .
\]
\end{theorem}

\begin{proof}
The first assertion follows from the polynomial dependence of
\(
        B(\Phi,\Phi)
\)
on the Higgs field, together with the Serre-duality identification
\[
        H^0(C,K_C^2(D))
        \cong
        T^*_{(C,D)}\mathcal T_{g,r}.
\]
Thus \(\phi_{G,\log}\) varies holomorphically along the smooth relative moduli
space. If \(\nu\) vanishes near \(D\), then the product \(B(\Phi,\Phi)\nu\) is a smooth
\((1,1)\)-form on \(C\).  In this case the Serre pairing is represented by the integral
\[
        \left\langle B(\Phi,\Phi),Y\right\rangle
        =
        \int_C B(\Phi,\Phi)\nu,
\]
which gives the required formula.
\end{proof}


\begin{thebibliography}{99}
\bibitem[BiswasGastesiGovindarajan]{BiswasGastesiGovindarajan} I.~Biswas, P.~Ar\'es-Gastesi and S.~Govindarajan, \emph{Parabolic Higgs bundles and Teichmuller spaces for punctured surfaces}, Trans. Amer. Math. Soc. \textbf{349} (1997), no.~4, 1551--1560.
\bibitem[BiswasParabolic]{BiswasParabolic} I.~Biswas, \emph{Parabolic bundles as orbifold bundles}, Duke Math. J. \textbf{88} (1997), no.~2, 305--325.
\bibitem[BodenYokogawa]{BodenYokogawa} H.~U. Boden and K.~Yokogawa, \emph{Moduli spaces of parabolic Higgs bundles and parabolic \(K(D)\)-pairs over smooth curves. I}, Internat. J. Math. \textbf{7} (1996), no.~5, 573--598.
\bibitem[CollingwoodMcGovern]{CollingwoodMcGovern} D.~H. Collingwood and W.~M. McGovern, \emph{Nilpotent Orbits in Semisimple Lie Algebras}, Van Nostrand Reinhold, New York, 1993.
\bibitem[Gardiner]{Gardiner} F.~P. Gardiner, \emph{Teichmuller Theory and Quadratic Differentials}, Wiley-Interscience, New York, 1987.
\bibitem[GriffithsHarris]{GriffithsHarris} P.~Griffiths and J.~Harris, \emph{Principles of Algebraic Geometry}, Wiley-Interscience, New York, 1978.
\bibitem[Hubbard]{Hubbard} J.~H. Hubbard, \emph{Teichmuller Theory and Applications to Geometry, Topology, and Dynamics. Vol. 1}, Matrix Editions, Ithaca, NY, 2006.
\bibitem[MehtaSeshadri]{MehtaSeshadri} V.~B. Mehta and C.~S. Seshadri, \emph{Moduli of vector bundles on curves with parabolic structures}, Math. Ann. \textbf{248} (1980), no.~3, 205--239.
\bibitem[Yokogawa]{Yokogawa} K.~Yokogawa, \emph{Compactification of moduli of parabolic sheaves and moduli of parabolic Higgs sheaves}, J. Math. Kyoto Univ. \textbf{33} (1993), no.~2, 451--504.
\end{thebibliography}
\end{document}
